\chapter{Conclusiones y trabajo futuro}
\label{ch:conclusiones}
% Fuentes futuras: objetivos aprobados; resultados canónicos; TFM_CLOSEOUT.md §5.
\section{Conclusiones respecto a los objetivos}
\pendiente{Redactar tras revisar Resultados y Discusión. Responder a cada
objetivo: realización técnica, aprendizaje experimental, utilidad y problemas
pendientes. No introducir cifras ni resultados nuevos.}
\section{Trabajo futuro}
\subsection{Operación}
\pendiente{Ingestión diaria, planificación de ejecuciones y actualización
acumulativa; distinguir fases CLI implementadas de automatización pendiente.}
\pendiente{FIGURA F09 — Arquitectura propuesta para una futura operación
diaria.\par Etiqueta prevista: \path{fig:operacion-diaria-futura}.
Título y leyenda deben declarar
TRABAJO FUTURO. Identificar los mecanismos actuales: CLI por etapas,
snapshots, revisión, trazabilidad y downstream regenerable. Separar las
propuestas: scheduler, coordinación automática, acumulación incremental,
reportes persistentes y cola unificada. Flujo previsto: BOE diario, extracción/revisión,
resolución territorial, grouping, Gold y Streamlit. Propiedades deseadas:
incremental, idempotente, trazable, reanudable, no bloqueante y regenerable;
no son garantías del sistema actual completo. El diseño deberá conservar el
último Gold validado mientras se revisan nuevas incidencias, sin eludir los
controles bloqueantes existentes. Referenciar desde esta subsección.}
\subsection{Calidad y gobierno}
\subsubsection{Cola unificada de incidencias y reportes de usuarios}
\label{sec:cola-unificada-futura}
El sistema incorpora detección y revisión humana para algunos tipos de
incidencias, pero no dispone de una cola persistente que reúna extracción,
territorio y reportes de usuarios. La extracción y la salvaguarda P0
comparten una cola; la resolución territorial conserva estados propios;
el control público de reporte abre un correo sin almacenarlo en la aplicación.

Como evolución posterior al TFM se propone registrar de forma persistente
las incidencias automáticas y los reportes de Streamlit en una cola común,
con identificador, origen, estado e historial de revisión. Una administración
de estos casos permitiría asignar y documentar su resolución. La aceptación
de un reporte no modificaría directamente Gold: una decisión que exigiera
cambiar datos debería producir una corrección versionada, seguida de
regeneración y validación antes de publicar una nueva salida.

Esta propuesta reutilizaría los mecanismos actuales de trazabilidad,
decisiones y materialización, y añadiría la recepción persistente de reportes,
su integración y la coordinación administrativa que aún no existen.
Debería conservar el vínculo entre incidencia, evidencia, decisión y versión
publicada, además del último Gold validado mientras haya casos bloqueantes.
Se integrará en la figura prevista de operación diaria futura, sin mezclarla
con el flujo implementado de la sección~\ref{sec:revision-humana}.

\pendiente{Mejoras motivadas por evaluación: alias, granularidad de eventos,
evidencia y propagación jerárquica, sin ajustar retrospectivamente el experimento.}
\pendiente{Valorar municipio y promotor como corroboración futura de identidad,
con evidencia y validación propias. Ninguno participa en la clave actual;
no trasladarlos retrospectivamente a la figura del grouping implementado.}
\subsection{Enriquecimiento funcional}
\pendiente{Promotores, potencia instalada o autorizada, almacenamiento,
hibridaciones, resultado ambiental, siguiente hito y relaciones adicionales,
solo como trabajo futuro. No denominarlos producción energética.}
